%%%%%%%%%%%%%%%%%%%%%%%%%%%%%%%%%%%%%%%%%

%
% Important note:
% This template requires the resume.cls file to be in the same directory as the
% .tex file. The resume.cls file provides the resume style used for structuring the
% document.

%
% Creator Peilin Li
% Contact me via twitter/wechat: @pe1l1nl1
% linkedin.com/peill and/or github/ppeill
% Inspired by Peppa Pig 
%%%%%%%%%%%%%%%%%%%%%%%%%%%%%%%%%%%%%%%%

%----------------------------------------------------------------------------------------
%	PACKAGES AND OTHER DOCUMENT CONFIGURATIONS
%----------------------------------------------------------------------------------------

\documentclass{resume} % Use the custom resume.cls style

\usepackage[left=0.40in,top=0.3in,right=0.75in,bottom=0.1in]{geometry} % Document margins

% \usepackage{fontspec}
% \setmainfont{Times New Roman}

\usepackage[T1]{fontenc}  

\usepackage{fontawesome5}
% \usepackage{times}
\usepackage{multicol}
\newcommand{\tab}[1]{\hspace{.2667\textwidth}\rlap{#1}}
\newcommand{\itab}[1]{\hspace{0em}\rlap{#1}}
\setlength{\columnsep}{4cm}
\addtolength{\topmargin}{0.5in}


\name{Edward Varvak} 


% \address{Department of Mathematics \\ 5734 S University Ave \\ Chicago, IL 60637}
\address{\texttt{evarvak@gmail.com} \ | \ \texttt{https://homepages.math.uic.edu/{\textasciitilde}evarvak/}}

\begin{document}

\begin{rSection}{Profile}
Analytical problem solver and machine learning practitioner with a Ph.D.-level
mathematical background. Expertise in deconstructing complex problems and developing
cutting-edge data-driven solutions. Proven skills in Python, computer vision, and deep
learning, demonstrated through successful industry internships. Excellent communicator
skilled at translating technical findings for diverse audiences.
\end{rSection}


\begin{rSection}{Work Experience}

{\bf Innovation Intern at PTC}, {\it Summer 2021} \\ 
Researched and prototyped applications for state-of-the-art deep learning models for
client use cases, focusing on computer vision and dimensionality reduction. 
\begin{itemize} \itemsep -0.5em \vspace{-0.5em} 
    \item Investigated and benchmarked automated object detection models to identify
    optimal architectures for resource-constrained environments.
    \item Co-developed an algorithm for estimating process cycle times from 
    video footage, reducing the need for manual video analysis.
\end{itemize} 

{\bf Machine Learning Intern at CIMCON}, {\it Summer 2019} \\
Developed and deployed a real-time object-detection prototype on the NVIDIA Jetson 
Nano, demonstrating a proof-of-concept for a new product line.
\begin{itemize} \itemsep -0.5em \vspace{-0.5em}
    \item Engineered a gun detection model using TensorFlow and Keras for real-time
    video analysis.
    \item Evaluated low-power hardware, optimizing the model for continuous edge 
    deployment.
    \item Presented technical results and project progress to R\&D, sales teams, and senior leadership.
\end{itemize}

{\bf Computer Vision Software Intern at Eonite}, {\it Summer 2017} \\
Contributed to an R\&D team developing software for 3D mapping and localization in 
indoor environments.
\begin{itemize} \itemsep -0.5em \vspace{-0.5em} 
    \item Improved camera-calibration tooling to reduce localization error and 
    increase mapping accuracy. 
    \item Implemented a C++ bundle-adjustment module to improve frame-to-frame 
    localization by minimizing reprojection error.
\end{itemize}

\end{rSection}




\begin{rSection}{Education}

{\bf Doctor of Philosophy (Ph.D.) in Mathematics} \\
{\it University of Chicago} \hfill {\it Sep 2026--May 2027 (transfer)} \\
{\it University of Illinois Chicago} \hfill {\it Aug 2022--Sep 2026}
\begin{itemize} \itemsep -0.5em \vspace{-0.5em}
    \item Repeat recipient of Research Training Group Fellowship, 4.0/4.0 GPA.
    \item Produced two first-author publications, invited speaker at numerous seminars 
    and conferences.
\end{itemize}

{\bf Bachelor of Science (B.S.) in Applied Mathematics and Statistics} \\
{\it Stony Brook University} \hfill {\it Sep 2018--May 2022}
\begin{itemize} \itemsep -0.5em \vspace{-0.5em}
    \item Honors Track in Mathematics, Cum Laude, 3.6/4.0 GPA.
    \item Recipient of the Kuga-Sah Award for Excellence in Mathematics.
\end{itemize}

\end{rSection}


\begin{rSection}{Technical Skills} 
{\bf Programming:} Python (proficient), C++, C \\
{\bf ML Frameworks:} PyTorch, Scikit-learn, TensorFlow \\
{\bf Tools and Hardware:} Git, Debian Linux, Docker, NVIDIA Jetson \\
{\bf Core Competencies:} Mathematical Modeling, Data Analysis, Problem Solving, Technical Presentations

\end{rSection}



% \begin{rSection}{Research and Publications}
% \begin{enumerate}
% \item \textit{The Mumford--Tate group of a very general Lagrangian fiber}, 
% TBD*, 2026

% \item \textit{The Monodromy of Compact Lagrangian Fibrations}, 16 pages, 
% {\bf IMRN}, rnag057, 2026

% \item \textit{Recovering Plane Nodal Curves from their Bitangents}, 12 pages, 
% {\bf Undergraduate Honors Thesis}, \\
% advised by Olivier Martin, 2022.

% \end{enumerate}

% \end{rSection}





% \begin{rSection}{Organization and Service}
% {Co-organized Graduate Algebraic Geometry Seminar (GAGS)} \hfill \textit{Fall 2025-Spring 2026} \\
% {Judged for QED Compitition (Math Circles of Chicago)} \hfill \textit{December 2025}\\
% {Organized Weil Conjecture Reading Group} \hfill \textit{Fall 2025} \\
% {Organized Lie Groups Seminar} \hfill \textit{Fall 2024}
% \end{rSection}



% \begin{rSection}{Honors and Awards}
% {Bernard Kurtin Award for research in mathematics, UIC MSCS Department} 
% \hfill \textit{2025-2026} \\
% {UIC RTG Grant} \hfill \textit{Fall 2025,} \\
% \mbox{} \hfill \textit{Summer 2025,} \\
% \mbox{} \hfill \textit{Spring-Summer 2024,} \\
% \mbox{} \hfill \textit{Spring-Summer 2023} \\
% {Honors in Mathematics, SBU Mathematics Department} \hfill {\em 2018-2022} \\
% {Kuga-Sah Award for excellence in mathematics, SBU Mathematics Department} 
% \hfill {\em 2022} \\
% {Cum Laude, Stony Brook University} \hfill {\em 2022}
% \end{rSection}




% \begin{rSection}{Teaching Experience}
% \textbf{University of Chicago}: \\
% Math 15200: Calculus II (Graduate Student Lecturer) \hfill \textit{Fall 2026}\\
% \\
% \textbf{University of Illinois Chicago}: \\
% Math 181: Calculus II (Teaching Assistant) \hfill \textit{Spring 2026} \\
% Math 160: Finite Math for Business (Teaching Assistant) \hfill \textit{Spring 2025,} 
% \newline \mbox{} \hfill \textit{Fall 2024,} 
% \newline \mbox{} \hfill \textit{Fall 2023,} 
% \newline \mbox{} \hfill \textit{Fall 2022}
% \end{rSection}


% \begin{rSection}{Languages}

% English (Native) \\
% Russian (Fluent)

% \end{rSection}

% \begin{rSection}{Skills}
% \LaTeX, Python, C, R
% \end{rSection} 


\end{document}